%-----------------------------------------------------------------------------------------
% Autor dieser Vorlage:
% Stefan Macke (http://fachinformatiker-anwendungsentwicklung.net)
% Permalink zur Vorlage: http://fiae.link/LaTeXVorlageFIAE
%
% Sämtliche verwendeten Abbildungen, Tabellen und Listings stammen von Dirk Grashorn.
%
% Lizenz: Creative Commons 4.0 Namensnennung - Weitergabe unter gleichen Bedingungen
% -----------------------------------------------------------------------------------------

\documentclass[
	ngerman,
	12pt,
	toc=listof, % Abbildungsverzeichnis sowie Tabellenverzeichnis in das Inhaltsverzeichnis aufnehmen
	toc=bibliography, % Literaturverzeichnis in das Inhaltsverzeichnis aufnehmen
	footnotes=multiple, % Trennen von direkt aufeinander folgenden Fußnoten
	parskip=half, % vertikalen Abstand zwischen Absätzen verwenden anstatt horizontale Einrückung von Folgeabsätzen
	numbers=noendperiod % Den letzten Punkt nach einer Nummerierung entfernen (nach DIN 5008)
]{scrartcl}
\pdfminorversion=5 % erlaubt das Einfügen von pdf-Dateien bis Version 1.7, ohne eine Fehlermeldung zu werfen (keine Garantie für fehlerfreies Einbetten!)
\usepackage[utf8]{inputenc} % muss als erstes eingebunden werden, da Meta/Packages ggfs. Sonderzeichen enthalten

% !TEX root = Projektdokumentation.tex

% Hinweis: der Titel muss zum Inhalt des Projekts passen und den zentralen Inhalt des Projekts deutlich herausstellen
\newcommand{\titel}{Entwicklung eines Kalkulationskatalogs}
\newcommand{\untertitel}{Webanwendung zur Verwaltung von Kalkulationspositionen}
\newcommand{\kompletterTitel}{\titel{} -- \untertitel}

\newcommand{\autorName}{Norman Paschke}
\newcommand{\autorAnschrift}{Am Hirschgrund 26}
\newcommand{\autorOrt}{09385 Lugau}

\newcommand{\betriebLogo}{DMK_Logo_print.pdf}
\newcommand{\betriebName}{DMK E-BUSINESS GmbH}
\newcommand{\betriebAnschrift}{Dresdner Straße 40}
\newcommand{\betriebOrt}{09130 Chemnitz}

\newcommand{\ausbildungsberuf}{Fachinformatiker für Anwendungsentwicklung}
\newcommand{\betreff}{Dokumentation zur betrieblichen Projektarbeit}
\newcommand{\pruefungstermin}{Sommer 2016}
\newcommand{\abgabeOrt}{Chemnitz}
\newcommand{\abgabeTermin}{30.05.2016}
 % Metadaten zu diesem Dokument (Autor usw.)
% !TEX root = ../Projektdokumentation.tex

% Anpassung an Landessprache ---------------------------------------------------
\usepackage{babel}

% Umlaute ----------------------------------------------------------------------
%   Umlaute/Sonderzeichen wie äüöß direkt im Quelltext verwenden (CodePage).
%   Erlaubt automatische Trennung von Worten mit Umlauten.
% ------------------------------------------------------------------------------
\usepackage[T1]{fontenc}
\usepackage{textcomp} % Euro-Zeichen etc.

% Schrift ----------------------------------------------------------------------
\usepackage{lmodern} % bessere Fonts
\usepackage{relsize} % Schriftgröße relativ festlegen

% Tabellen ---------------------------------------------------------------------
\PassOptionsToPackage{table}{xcolor}
\usepackage{tabularx}
% für lange Tabellen
\usepackage{longtable}
\usepackage{array}
\usepackage{ragged2e}
\usepackage{lscape}
\newcolumntype{w}[1]{>{\raggedleft\hspace{0pt}}p{#1}} % Spaltendefinition rechtsbündig mit definierter Breite

% Grafiken ---------------------------------------------------------------------
\usepackage[dvips,final]{graphicx} % Einbinden von JPG-Grafiken ermöglichen
\usepackage{graphics} % keepaspectratio
\usepackage{floatflt} % zum Umfließen von Bildern
\graphicspath{{Bilder/}} % hier liegen die Bilder des Dokuments

% Sonstiges --------------------------------------------------------------------
\usepackage[titles]{tocloft} % Inhaltsverzeichnis DIN 5008 gerecht einrücken
\usepackage{amsmath,amsfonts} % Befehle aus AMSTeX für mathematische Symbole
\usepackage{enumitem} % anpassbare Enumerates/Itemizes
\usepackage{xspace} % sorgt dafür, dass Leerzeichen hinter parameterlosen Makros nicht als Makroendezeichen interpretiert werden

\usepackage{makeidx} % für Index-Ausgabe mit \printindex
\usepackage[printonlyused]{acronym} % es werden nur benutzte Definitionen aufgelistet

% Einfache Definition der Zeilenabstände und Seitenränder etc.
\usepackage{setspace}
\usepackage{geometry}

% Symbolverzeichnis
\usepackage[intoc]{nomencl}
\let\abbrev\nomenclature
\renewcommand{\nomname}{Abkürzungsverzeichnis}
\setlength{\nomlabelwidth}{.25\hsize}
\renewcommand{\nomlabel}[1]{#1 \dotfill}
\setlength{\nomitemsep}{-\parsep}

\usepackage{varioref} % Elegantere Verweise. „auf der nächsten Seite“
\usepackage{url} % URL verlinken, lange URLs umbrechen etc.

\usepackage{chngcntr} % fortlaufendes Durchnummerieren der Fußnoten
% \usepackage[perpage]{footmisc} % Alternative: Nummerierung der Fußnoten auf jeder Seite neu

\usepackage{ifthen} % bei der Definition eigener Befehle benötigt
\usepackage{todonotes} % definiert u.a. die Befehle \todo und \listoftodos
\usepackage[square]{natbib} % wichtig für korrekte Zitierweise

% PDF-Optionen -----------------------------------------------------------------
\usepackage{pdfpages}
\pdfminorversion=5 % erlaubt das Einfügen von pdf-Dateien bis Version 1.7, ohne eine Fehlermeldung zu werfen (keine Garantie für fehlerfreies Einbetten!)
\usepackage[
    bookmarks,
    bookmarksnumbered,
    bookmarksopen=true,
    bookmarksopenlevel=1,
    colorlinks=true,
% diese Farbdefinitionen zeichnen Links im PDF farblich aus
    linkcolor=AOBlau, % einfache interne Verknüpfungen
    anchorcolor=AOBlau,% Ankertext
    citecolor=AOBlau, % Verweise auf Literaturverzeichniseinträge im Text
    filecolor=AOBlau, % Verknüpfungen, die lokale Dateien öffnen
    menucolor=AOBlau, % Acrobat-Menüpunkte
    urlcolor=AOBlau,
% diese Farbdefinitionen sollten für den Druck verwendet werden (alles schwarz)
    %linkcolor=black, % einfache interne Verknüpfungen
    %anchorcolor=black, % Ankertext
    %citecolor=black, % Verweise auf Literaturverzeichniseinträge im Text
    %filecolor=black, % Verknüpfungen, die lokale Dateien öffnen
    %menucolor=black, % Acrobat-Menüpunkte
    %urlcolor=black,
%
    %backref, % Quellen werden zurück auf ihre Zitate verlinkt
    pdftex,
    plainpages=false, % zur korrekten Erstellung der Bookmarks
    pdfpagelabels=true, % zur korrekten Erstellung der Bookmarks
    hypertexnames=false, % zur korrekten Erstellung der Bookmarks
    linktocpage % Seitenzahlen anstatt Text im Inhaltsverzeichnis verlinken
]{hyperref}
% Befehle, die Umlaute ausgeben, führen zu Fehlern, wenn sie hyperref als Optionen übergeben werden
\hypersetup{
    pdftitle={\titel -- \untertitel},
    pdfauthor={\autorName},
    pdfcreator={\autorName},
    pdfsubject={\titel -- \untertitel},
    pdfkeywords={\titel -- \untertitel},
}


% zum Einbinden von Programmcode -----------------------------------------------
\usepackage{listings}
\usepackage{xcolor}
\definecolor{hellgelb}{rgb}{1,1,0.9}
\definecolor{colKeys}{rgb}{0,0,1}
\definecolor{colIdentifier}{rgb}{0,0,0}
\definecolor{colComments}{rgb}{0,0.5,0}
\definecolor{colString}{rgb}{1,0,0}
\lstset{
    float=hbp,
	basicstyle=\footnotesize,
    identifierstyle=\color{colIdentifier},
    keywordstyle=\color{colKeys},
    stringstyle=\color{colString},
    commentstyle=\color{colComments},
    backgroundcolor=\color{hellgelb},
    columns=flexible,
    tabsize=2,
    frame=single,
    extendedchars=true,
    showspaces=false,
    showstringspaces=false,
    numbers=left,
    numberstyle=\tiny,
    breaklines=true,
    breakautoindent=true,
	captionpos=b,
}
\lstdefinelanguage{cs}{
	sensitive=false,
	morecomment=[l]{//},
	morecomment=[s]{/*}{*/},
	morestring=[b]",
	morekeywords={
		abstract,event,new,struct,as,explicit,null,switch
		base,extern,object,this,bool,false,operator,throw,
		break,finally,out,true,byte,fixed,override,try,
		case,float,params,typeof,catch,for,private,uint,
		char,foreach,protected,ulong,checked,goto,public,unchecked,
		class,if,readonly,unsafe,const,implicit,ref,ushort,
		continue,in,return,using,decimal,int,sbyte,virtual,
		default,interface,sealed,volatile,delegate,internal,short,void,
		do,is,sizeof,while,double,lock,stackalloc,
		else,long,static,enum,namespace,string},
}
\lstdefinelanguage{natural}{
	sensitive=false,
	morecomment=[l]{/*},
	morestring=[b]",
	morestring=[b]',
	alsodigit={-,*},
	morekeywords={
		DEFINE,DATA,LOCAL,END-DEFINE,WRITE,CALLNAT,PARAMETER,USING,
		IF,NOT,END-IF,ON,*ERROR-NR,ERROR,END-ERROR,ESCAPE,ROUTINE,
		PERFORM,SUBROUTINE,END-SUBROUTINE,CONST,END-FOR,END,FOR,RESIZE,
		ARRAY,TO,BY,VALUE,RESET,COMPRESS,INTO,EQ},
}
\lstdefinelanguage{php}{
	sensitive=false,
	morecomment=[l]{/*},
	morestring=[b]",
	morestring=[b]',
	alsodigit={-,*},
	morekeywords={
		abstract,and,array,as,break,case,catch,cfunction,class,clone,const,
		continue,declare,default,do,else,elseif,enddeclare,endfor,endforeach,
		endif,endswitch,endwhile,extends,final,for,foreach,function,global,
		goto,if,implements,interface,instanceof,namespace,new,old_function,or,
		private,protected,public,static,switch,throw,try,use,var,while,xor
		die,echo,empty,exit,eval,include,include_once,isset,list,require,
		require_once,return,print,unset},
}
 % verwendete Packages
% !TEX root = ../Projektdokumentation.tex

% Seitenränder -----------------------------------------------------------------
\setlength{\topskip}{\ht\strutbox} % behebt Warnung von geometry
\geometry{a4paper,left=20mm,right=20mm,top=25mm,bottom=35mm}

\usepackage[
	automark, % Kapitelangaben in Kopfzeile automatisch erstellen
	headsepline, % Trennlinie unter Kopfzeile
	ilines % Trennlinie linksbündig ausrichten
]{scrpage2}

% Kopf- und Fußzeilen ----------------------------------------------------------
\pagestyle{scrheadings}
% chapterpagestyle gibt es nicht in scrartcl
%\renewcommand{\chapterpagestyle}{scrheadings}
\clearscrheadfoot

% Kopfzeile
\renewcommand{\headfont}{\normalfont} % Schriftform der Kopfzeile
\ihead{\textit{\headmark}}
\chead{}
%\ohead{\includegraphics[scale=0.2]{\betriebLogo}}
\setlength{\headheight}{10mm} % Höhe der Kopfzeile
%\setheadwidth[0pt]{textwithmarginpar} % Kopfzeile über den Text hinaus
% verbreitern (falls Logo den Text überdeckt)

% Fußzeile
\cfoot{\pagemark}

% Überschriften nach DIN 5008 in einer Fluchtlinie
% ------------------------------------------------------------------------------

% Abstand zwischen Nummerierung und Überschrift definieren
% > Schön wäre hier die dynamische Berechnung des Abstandes in Abhängigkeit
% > der Verschachtelungstiefe des Inhaltsverzeichnisses
\newcommand{\headingSpace}{1.5cm}

% Abschnittsüberschriften im selben Stil wie beim Inhaltsverzeichnis einrücken
\renewcommand*{\othersectionlevelsformat}[3]{
  \makebox[\headingSpace][l]{#3\autodot}
}

% Für die Einrückung wird das Paket tocloft benötigt
%\cftsetindents{chapter}{0.0cm}{\headingSpace}
\cftsetindents{section}{0.0cm}{\headingSpace}
\cftsetindents{subsection}{0.0cm}{\headingSpace}
\cftsetindents{subsubsection}{0.0cm}{\headingSpace}
\cftsetindents{figure}{0.0cm}{\headingSpace}
\cftsetindents{table}{0.0cm}{\headingSpace}


% Allgemeines
% ------------------------------------------------------------------------------

\onehalfspacing % Zeilenabstand 1,5 Zeilen
\frenchspacing % erzeugt ein wenig mehr Platz hinter einem Punkt

% Schusterjungen und Hurenkinder vermeiden
\clubpenalty = 10000
\widowpenalty = 10000
\displaywidowpenalty = 10000

% Quellcode-Ausgabe formatieren
\lstset{numbers=left, numberstyle=\tiny, numbersep=5pt, breaklines=true}
\lstset{emph={square}, emphstyle=\color{red}, emph={[2]root,base}, emphstyle={[2]\color{blue}}}

\counterwithout{footnote}{section} % Fußnoten fortlaufend durchnummerieren
\setcounter{tocdepth}{\subsubsectionlevel} % im Inhaltsverzeichnis werden die Kapitel bis zum Level der subsubsection übernommen
\setcounter{secnumdepth}{\subsubsectionlevel} % Kapitel bis zum Level der subsubsection werden nummeriert

% Aufzählungen anpassen
\renewcommand{\labelenumi}{\arabic{enumi}.}
\renewcommand{\labelenumii}{\arabic{enumi}.\arabic{enumii}.}
\renewcommand{\labelenumiii}{\arabic{enumi}.\arabic{enumii}.\arabic{enumiii}}

% Tabellenfärbung:
\definecolor{heading}{rgb}{0.64,0.78,0.86}
\definecolor{odd}{rgb}{0.9,0.9,0.9}
 % Definitionen zum Aussehen der Seiten
% !TEX root = ../Projektdokumentation.tex

% Abkürzungen, ggfs. mit korrektem Leerraum
\newcommand{\bs}{$\backslash$\xspace}
\newcommand{\bspw}{bspw.\xspace}
\newcommand{\bzw}{bzw.\xspace}
\newcommand{\ca}{ca.\xspace}
\newcommand{\dahe}{\mbox{d.\,h.}\xspace}
\newcommand{\etc}{etc.\xspace}
\newcommand{\eur}[1]{\mbox{#1\,\texteuro}\xspace}
\newcommand{\evtl}{evtl.\xspace}
\newcommand{\ggfs}{ggfs.\xspace}
\newcommand{\Ggfs}{Ggfs.\xspace}
\newcommand{\gqq}[1]{\glqq{}#1\grqq{}}
\newcommand{\inkl}{inkl.\xspace}
\newcommand{\insb}{insb.\xspace}
\newcommand{\ua}{\mbox{u.\,a.}\xspace}
\newcommand{\usw}{usw.\xspace}
\newcommand{\Vgl}{Vgl.\xspace}
\newcommand{\zB}{\mbox{z.\,B.}\xspace}

% Befehle für häufig anfallende Aufgaben
\newcommand{\Abbildung}[1]{\autoref{fig:#1}}
\newcommand{\Anhang}[1]{\appendixname{}~\ref{#1}: \nameref{#1} \vpageref{#1}}
\newcommand{\includegraphicsKeepAspectRatio}[2]{\includegraphics[width=#2\textwidth,height=#2\textheight,keepaspectratio]{#1}}
\newcommand{\Zitat}[2][\empty]{\ifthenelse{\equal{#1}{\empty}}{\citep{#2}}{\citep[#1]{#2}}}
\newcommand{\Autor}[1]{\textsc{#1}} % zum Ausgeben von Autoren
\newcommand{\itemd}[2]{\item{\textbf{#1}}\\{#2}} % erzeugt ein Listenelement mit fetter Überschrift

% fügt Tabellen aus einer TEX-Datei ein
\newcommand{\tabelle}[3] % Parameter: caption, label, file
{\begin{table}[htbp]
\centering
\singlespacing
\input{Tabellen/#3}
\caption{#1}
\label{#2}
\end{table}}

\newcommand{\tabelleAnhang}[1] % Parameter: file
{\begin{center}
\singlespacing
\input{Tabellen/#1}
\end{center}}

% einfaches Wechseln der Schrift, z.B.: \changefont{cmss}{sbc}{n}
\newcommand{\changefont}[3]{\fontfamily{#1} \fontseries{#2} \fontshape{#3} \selectfont}

% Verwendung analog zu \includegraphics
\newlength{\myx} % Variable zum Speichern der Bildbreite
\newlength{\myy} % Variable zum Speichern der Bildhöhe
\newcommand\includegraphicstotab[2][\relax]{%
% Abspeichern der Bildabmessungen
\settowidth{\myx}{\includegraphics[{#1}]{#2}}%
\settoheight{\myy}{\includegraphics[{#1}]{#2}}%
% das eigentliche Einfügen
\parbox[c][1.1\myy][c]{\myx}{%
\includegraphics[{#1}]{#2}}%
}

\definecolor{AOBlau}{rgb}{0, 0.28, 0.56}

% verschiedene Befehle um Wörter semantisch auszuzeichnen ----------------------
\newcommand{\Index}[2][\empty]{\ifthenelse{\equal{#1}{\empty}}{\index{#2}#2}{\index{#1}#2}}
\newcommand{\Fachbegriff}[2][\empty]{\ifthenelse{\equal{#1}{\empty}}{\textit{\Index{#2}}}{\textit{\Index[#1]{#2}}}}
\newcommand{\NeuerBegriff}[2][\empty]{\ifthenelse{\equal{#1}{\empty}}{\textbf{\Index{#2}}}{\textbf{\Index[#1]{#2}}}}

\newcommand{\Ausgabe}[1]{\texttt{#1}}
\newcommand{\Eingabe}[1]{\texttt{#1}}
\newcommand{\Code}[1]{\texttt{#1}}
\newcommand{\Datei}[1]{\texttt{#1}}

\newcommand{\Assembly}[1]{\textsf{#1}}
\newcommand{\Klasse}[1]{\textsf{#1}}
\newcommand{\Methode}[1]{\textsf{#1}}
\newcommand{\Attribut}[1]{\textsf{#1}}

\newcommand{\Datentyp}[1]{\textsf{#1}}
\newcommand{\XMLElement}[1]{\textsf{#1}}
\newcommand{\Webservice}[1]{\textsf{#1}}

\newcommand{\Refactoring}[1]{\Fachbegriff{#1}}
\newcommand{\CodeSmell}[1]{\Fachbegriff{#1}}
\newcommand{\Metrik}[1]{\Fachbegriff{#1}}
\newcommand{\DesignPattern}[1]{\Fachbegriff{#1}}
 % eigene allgemeine Befehle, die z.B. die Arbeit mit LaTeX erleichtern
% !TEX root = ../Projektdokumentation.tex

% Abkürzungen, ggfs. mit korrektem Leerraum
\newcommand{\bs}{$\backslash$\xspace}
\newcommand{\bspw}{bspw.\xspace}
\newcommand{\bzw}{bzw.\xspace}
\newcommand{\ca}{ca.\xspace}
\newcommand{\dahe}{\mbox{d.\,h.}\xspace}
\newcommand{\etc}{etc.\xspace}
\newcommand{\eur}[1]{\mbox{#1\,\texteuro}\xspace}
\newcommand{\evtl}{evtl.\xspace}
\newcommand{\ggfs}{ggfs.\xspace}
\newcommand{\Ggfs}{Ggfs.\xspace}
\newcommand{\gqq}[1]{\glqq{}#1\grqq{}}
\newcommand{\inkl}{inkl.\xspace}
\newcommand{\insb}{insb.\xspace}
\newcommand{\ua}{\mbox{u.\,a.}\xspace}
\newcommand{\usw}{usw.\xspace}
\newcommand{\Vgl}{Vgl.\xspace}
\newcommand{\zB}{\mbox{z.\,B.}\xspace}

% Befehle für häufig anfallende Aufgaben
\newcommand{\Abbildung}[1]{\autoref{fig:#1}}
\newcommand{\Anhang}[1]{\appendixname{}~\ref{#1}: \nameref{#1} \vpageref{#1}}
\newcommand{\includegraphicsKeepAspectRatio}[2]{\includegraphics[width=#2\textwidth,height=#2\textheight,keepaspectratio]{#1}}
\newcommand{\Zitat}[2][\empty]{\ifthenelse{\equal{#1}{\empty}}{\citep{#2}}{\citep[#1]{#2}}}
\newcommand{\Autor}[1]{\textsc{#1}} % zum Ausgeben von Autoren
\newcommand{\itemd}[2]{\item{\textbf{#1}}\\{#2}} % erzeugt ein Listenelement mit fetter Überschrift

% fügt Tabellen aus einer TEX-Datei ein
\newcommand{\tabelle}[3] % Parameter: caption, label, file
{\begin{table}[htbp]
\centering
\singlespacing
\input{Tabellen/#3}
\caption{#1}
\label{#2}
\end{table}}

\newcommand{\tabelleAnhang}[1] % Parameter: file
{\begin{center}
\singlespacing
\input{Tabellen/#1}
\end{center}}

% einfaches Wechseln der Schrift, z.B.: \changefont{cmss}{sbc}{n}
\newcommand{\changefont}[3]{\fontfamily{#1} \fontseries{#2} \fontshape{#3} \selectfont}

% Verwendung analog zu \includegraphics
\newlength{\myx} % Variable zum Speichern der Bildbreite
\newlength{\myy} % Variable zum Speichern der Bildhöhe
\newcommand\includegraphicstotab[2][\relax]{%
% Abspeichern der Bildabmessungen
\settowidth{\myx}{\includegraphics[{#1}]{#2}}%
\settoheight{\myy}{\includegraphics[{#1}]{#2}}%
% das eigentliche Einfügen
\parbox[c][1.1\myy][c]{\myx}{%
\includegraphics[{#1}]{#2}}%
}

\definecolor{AOBlau}{rgb}{0, 0.28, 0.56}

% verschiedene Befehle um Wörter semantisch auszuzeichnen ----------------------
\newcommand{\Index}[2][\empty]{\ifthenelse{\equal{#1}{\empty}}{\index{#2}#2}{\index{#1}#2}}
\newcommand{\Fachbegriff}[2][\empty]{\ifthenelse{\equal{#1}{\empty}}{\textit{\Index{#2}}}{\textit{\Index[#1]{#2}}}}
\newcommand{\NeuerBegriff}[2][\empty]{\ifthenelse{\equal{#1}{\empty}}{\textbf{\Index{#2}}}{\textbf{\Index[#1]{#2}}}}

\newcommand{\Ausgabe}[1]{\texttt{#1}}
\newcommand{\Eingabe}[1]{\texttt{#1}}
\newcommand{\Code}[1]{\texttt{#1}}
\newcommand{\Datei}[1]{\texttt{#1}}

\newcommand{\Assembly}[1]{\textsf{#1}}
\newcommand{\Klasse}[1]{\textsf{#1}}
\newcommand{\Methode}[1]{\textsf{#1}}
\newcommand{\Attribut}[1]{\textsf{#1}}

\newcommand{\Datentyp}[1]{\textsf{#1}}
\newcommand{\XMLElement}[1]{\textsf{#1}}
\newcommand{\Webservice}[1]{\textsf{#1}}

\newcommand{\Refactoring}[1]{\Fachbegriff{#1}}
\newcommand{\CodeSmell}[1]{\Fachbegriff{#1}}
\newcommand{\Metrik}[1]{\Fachbegriff{#1}}
\newcommand{\DesignPattern}[1]{\Fachbegriff{#1}}
 % eigene projektspezifische Befehle, z.B. Abkürzungen usw.

\begin{document}

\phantomsection
\thispagestyle{plain}
\pdfbookmark[1]{Deckblatt}{deckblatt}
% !TEX root = Projektdokumentation.tex
\begin{titlepage}

\begin{center}
\includegraphics[scale=0.25]{LogoIHK.pdf}\\[1ex]
\Large{Abschlussprüfung \pruefungstermin}\\[3ex]

\Large{\ausbildungsberuf}\\
\LARGE{\betreff}\\[4ex]

\huge{\textbf{\titel}}\\[1.5ex]
\Large{\textbf{\untertitel}}\\[4ex]

\normalsize
Abgabetermin: \abgabeOrt, den \abgabeTermin\\[5ex]
\textbf{Prüfungsbewerber:}\\
\autorName\\[2ex]

\includegraphics[scale=0.3]{\betriebLogo}\\[2ex]
\textbf{Ausbildungsbetrieb:}\\
\betriebName\\
\betriebAnschrift\\
\betriebOrt\\[2ex]
\textbf{Projektbetreuer:}\\
Tom Starke, B. Sc.
\end{center}

\end{titlepage}
\cleardoublepage

% Preface --------------------------------------------------------------------
\phantomsection
\pagenumbering{Roman}
\pdfbookmark[1]{Inhaltsverzeichnis}{inhalt}
\tableofcontents
\cleardoublepage

\phantomsection
\listoffigures
\cleardoublepage

\phantomsection
\listoftables
\cleardoublepage

%\phantomsection
%\lstlistoflistings
%\cleardoublepage

\newcommand{\abkvz}{Abkürzungsverzeichnis}
\renewcommand{\nomname}{\abkvz}
\section*{\abkvz}
\markboth{\abkvz}{\abkvz}
\addcontentsline{toc}{section}{\abkvz}
% !TEX root = Projektdokumentation.tex

% Es werden nur die Abkürzungen aufgelistet, die mit \ac definiert und auch benutzt wurden. 
%
% \acro{VERSIS}{Versicherungsinformationssystem\acroextra{ (Bestandsführungssystem)}}
% Ergibt in der Liste: VERSIS Versicherungsinformationssystem (Bestandsführungssystem)
% Im Text aber: \ac{VERSIS} -> Versicherungsinformationssystem (VERSIS)

% Hinweis: allgemein bekannte Abkürzungen wie z.B. bzw. u.a. müssen nicht ins Abkürzungsverzeichnis aufgenommen werden
% Hinweis: allgemein bekannte IT-Begriffe wie Datenbank oder Programmiersprache müssen nicht erläutert werden,
%          aber ggfs. Fachbegriffe aus der Domäne des Prüflings (z.B. Versicherung)

% Die Option (in den eckigen Klammern) enthält das längste Label oder
% einen Platzhalter der die Breite der linken Spalte bestimmt.
\begin{acronym}[WWWWW]
	\acro{API}{Application Programming Interface}
	\acro{BE}{Back-End}
	\acro{CMS}{Content Management System}
	\acro{CSS}{Cascading Style Sheets}
	\acro{CSV}{Comma Separated Value}
	\acro{DMK}{DMK E-BUSINESS GmbH}
	\acro{EPK}{Ereignisgesteuerte Prozesskette}
	\acro{ERM}{Entity-Relationship-Modell}
	\acro{FE}{Front-End}
	\acro{HTML}{Hypertext Markup Language}
	\acro{IDE}{Integrated Development Environment}
	\acro{MVC}{Model View Controller}
	\acro{PHP}{Hypertext Preprocessor}
	\acro{PM}{Projektmanager}
	\acro{SQL}{Structured Query Language}
	\acro{UML}{Unified Modeling Language}
\end{acronym}

\clearpage

% Inhalt ---------------------------------------------------------------------
\pagenumbering{arabic}
% !TEX root = Projektdokumentation.tex
% !TEX root = ../Projektdokumentation.tex
\section{Einleitung}
\label{sec:Einleitung}


\subsection{Projektumfeld} 
\label{sec:Projektumfeld}
\begin{itemize}
	\item Kurze Vorstellung des Ausbildungsbetriebs (Geschäftsfeld, Mitarbeiterzahl \usw)
	\item Wer ist Auftraggeber/Kunde des Projekts?
\end{itemize}


\subsection{Projektziel} 
\label{sec:Projektziel}
\begin{itemize}
	\item Worum geht es eigentlich?
	\item Was soll erreicht werden?
\end{itemize}


\subsection{Projektbegründung} 
\label{sec:Projektbegruendung}
\begin{itemize}
	\item Warum ist das Projekt sinnvoll (\zB Kosten- oder Zeitersparnis, weniger Fehler)?
	\item Was ist die Motivation hinter dem Projekt?
\end{itemize}


\subsection{Projektschnittstellen} 
\label{sec:Projektschnittstellen}
\begin{itemize}
	\item Mit welchen anderen Systemen interagiert die Anwendung (technische Schnittstellen)?
	\item Wer genehmigt das Projekt \bzw stellt Mittel zur Verfügung? 
	\item Wer sind die Benutzer der Anwendung?
	\item Wem muss das Ergebnis präsentiert werden?
\end{itemize}


\subsection{Projektabgrenzung} 
\label{sec:Projektabgrenzung}
\begin{itemize}
	\item Was ist explizit nicht Teil des Projekts (\insb bei Teilprojekten)?
\end{itemize}

% !TEX root = ../Projektdokumentation.tex
\section{Projektplanung} 
\label{sec:Projektplanung}

\subsection{Projektphasen}
\label{sec:Projektphasen}
Eine grobe Zeitplanung wurde im Vorfeld erstellt, dabei war sich an der Vorgabe
von 70~Stunden zu halten. Tabelle~\ref{tab:Zeitplanung}:
\nameref{tab:Zeitplanung} stellt den groben zeitlichen Ablauf des Projektes dar.
Die Entwicklung des Katalogs begann Mitte April 2016 und zog sich über sieben
Wochen hin, in denen regelmäßig ca. 12~Stunden pro Woche am Projekt gearbeitet
wurde.

\tabelle{Zeitplanung}{tab:Zeitplanung}{ZeitplanungKurz}

Die detaillierte Zeitplanung findet sich im \Anhang{app:Zeitplanung}.

\subsection{Ressourcenplanung}
\label{sec:Ressourcenplanung}
Alle im Projekt verwendeten Ressourcen finden sich in einer Übersicht im
\Anhang{app:Verwendete Ressourcen}. Benötigte Hardware wurde von \acs{DMK} zur
Verfügung gestellt. Bei der Auswahl der Software wurde darauf geachtet, dass die
Produkte kostenfrei (\zB Open Source) erhätlich und nutzbar sind oder -- wie im
Fall von Windows -- das Unternehmen bereits Lizenzen dafür besitzt.

\subsection{Entwicklungsprozess}
\label{sec:Entwicklungsprozess}
Bevor es in die Umsetzungsphase gehen konnte, war es nötig sich für einen
passenden Entwicklungsprozess zu entscheiden. Dabei galt es zwischen einem
Wasserfallmodell oder einem agilen Prozess abzuwägen. Um eine stetige
Rückmeldung von den Testern zu erhalten, fiel die Wahl schnell auf ein agiles
Vorgehen. Dabei wurde aber nicht explizit ein Vorgehen nach
Scrum\footnote{\url{https://de.wikipedia.org/wiki/Scrum}} favorisiert, sondern eine flexiblere Variante, die regelmäßige Treffen in Form von
Test-Sitzungen vorsah. Klassische Daily Stand-Ups oder Sprintplanungen fanden
daher nicht statt.

Die strikte Einteilung in Phasen, wie sie in~\ref{sec:Projektphasen}:
\nameref{sec:Projektphasen} vorgenommen wurde, musste entsprechend aufgeweicht
werden, da mehrere Phasen zyklisch auftraten und sich mit geringerem Aufwand
wiederholten.

Agiles Vorgehen hat den Vorteil, dass sehr schnell ein arbeitsfähiger Stand
präsentiert werden kann. Am Ende jeder Woche wurde eine kurze Vorstellung des
Projektes durchgeführt. Anmerkungen und Fehler wurden notiert und im nächsten
Zyklus bearbeitet, sodass eine kontinuierliche Verbesserung des Programms
stattgefunden hat.
% !TEX root = ../Projektdokumentation.tex
\section{Analysephase} 
\label{sec:Analysephase}


\subsection{Ist-Analyse} 
\label{sec:IstAnalyse}
\begin{itemize}
	\item Wie ist die bisherige Situation (\zB bestehende Programme, Wünsche der Mitarbeiter)?
	\item Was gilt es zu erstellen/verbessern?
\end{itemize}


\subsection{Wirtschaftlichkeitsanalyse}
\label{sec:Wirtschaftlichkeitsanalyse}
\begin{itemize}
	\item Lohnt sich das Projekt für das Unternehmen?
\end{itemize}


\subsubsection{\gqq{Make or Buy}-Entscheidung}
\label{sec:MakeOrBuyEntscheidung}
\begin{itemize}
	\item Gibt es vielleicht schon ein fertiges Produkt, dass alle Anforderungen des Projekts abdeckt?
	\item Wenn ja, wieso wird das Projekt trotzdem umgesetzt?
\end{itemize}


\subsubsection{Projektkosten}
\label{sec:Projektkosten}
\begin{itemize}
	\item Welche Kosten fallen bei der Umsetzung des Projekts im Detail an (\zB Entwicklung, Einführung/Schulung, Wartung)?
\end{itemize}

\paragraph{Beispielrechnung (verkürzt)}
Die Kosten für die Durchführung des Projekts setzen sich sowohl aus Personal-, als auch aus Ressourcenkosten zusammen.
Laut Tarifvertrag verdient ein Auszubildender im dritten Lehrjahr pro Monat \eur{1000} Brutto. 

\begin{eqnarray}
8 \mbox{ h/Tag} \cdot 220 \mbox{ Tage/Jahr} = 1760 \mbox{ h/Jahr}\\
\eur{1000}\mbox{/Monat} \cdot 13,3 \mbox{ Monate/Jahr} = \eur{13300} \mbox{/Jahr}\\
\frac{\eur{13300} \mbox{/Jahr}}{1760 \mbox{ h/Jahr}} \approx \eur{7,56}\mbox{/h}
\end{eqnarray}

Es ergibt sich also ein Stundenlohn von \eur{7,56}. 
Die Durchführungszeit des Projekts beträgt 70 Stunden. Für die Nutzung von Ressourcen\footnote{Räumlichkeiten, Arbeitsplatzrechner etc.} wird 
ein pauschaler Stundensatz von \eur{15} angenommen. Für die anderen Mitarbeiter wird pauschal ein Stundenlohn von \eur{25} angenommen. 
Eine Aufstellung der Kosten befindet sich in Tabelle~\ref{tab:Kostenaufstellung} und sie betragen insgesamt \eur{2739,20}.
\tabelle{Kostenaufstellung}{tab:Kostenaufstellung}{Kostenaufstellung.tex}


\subsubsection{Amortisationsdauer}
\label{sec:Amortisationsdauer}
\begin{itemize}
	\item Welche monetären Vorteile bietet das Projekt (\zB Einsparung von Lizenzkosten, Arbeitszeitersparnis, bessere Usability, Korrektheit)?
	\item Wann hat sich das Projekt amortisiert?
\end{itemize}

\paragraph{Beispielrechnung (verkürzt)}
Bei einer Zeiteinsparung von 10 Minuten am Tag für jeden der 25 Anwender und 220 Arbeitstagen im Jahr ergibt sich eine gesamte Zeiteinsparung von 
\begin{eqnarray}
25 \cdot 220 \mbox{ Tage/Jahr} \cdot 10 \mbox{ min/Tag} = 55000 \mbox{ min/Jahr} \approx 917 \mbox{ h/Jahr} 
\end{eqnarray}

Dadurch ergibt sich eine jährliche Einsparung von 
\begin{eqnarray}
917 \mbox{h} \cdot \eur{(25 + 15)}{\mbox{/h}} = \eur{36680}
\end{eqnarray}

Die Amortisationszeit beträgt also $\frac{\eur{2739,20}}{\eur{36680}\mbox{/Jahr}} \approx 0,07 \mbox{ Jahre} \approx 4 \mbox{ Wochen}$.


\subsection{Nutzwertanalyse}
\label{sec:Nutzwertanalyse}
\begin{itemize}
	\item Darstellung des nicht-monetären Nutzens (\zB Vorher-/Nachher-Vergleich anhand eines Wirtschaftlichkeitskoeffizienten). 
\end{itemize}

\paragraph{Beispiel}
Ein Beispiel für eine Entscheidungsmatrix findet sich in Kapitel~\ref{sec:Architekturdesign}: \nameref{sec:Architekturdesign}.


\subsection{Anwendungsfälle}
\label{sec:Anwendungsfaelle}
\begin{itemize}
	\item Welche Anwendungsfälle soll das Projekt abdecken?
	\item Einer oder mehrere interessante (!) Anwendungsfälle könnten exemplarisch durch ein Aktivitätsdiagramm oder eine \ac{EPK} detailliert beschrieben werden. 
\end{itemize}

\paragraph{Beispiel}
Ein Beispiel für ein Use Case-Diagramm findet sich im \Anhang{app:UseCase}.


\subsection{Qualitätsanforderungen}
\label{sec:Qualitaetsanforderungen}
\begin{itemize}
	\item Welche Qualitätsanforderungen werden an die Anwendung gestellt (\zB hinsichtlich Performance, Usability, Effizienz \etc (siehe \citet{ISO9126}))?
\end{itemize}


\subsection{Lastenheft/Fachkonzept}
\label{sec:Lastenheft}
\begin{itemize}
	\item Auszüge aus dem Lastenheft/Fachkonzept, wenn es im Rahmen des Projekts erstellt wurde.
	\item Mögliche Inhalte: Funktionen des Programms (Muss/Soll/Wunsch), User Stories, Benutzerrollen
\end{itemize}

\paragraph{Beispiel}
Ein Beispiel für ein Lastenheft findet sich im \Anhang{app:Lastenheft}. 

\Zwischenstand{Analysephase}{Analyse}

% !TEX root = ../Projektdokumentation.tex
\section{Entwurfsphase} 
\label{sec:Entwurfsphase}

\subsection{Zielplattform}
\label{sec:Zielplattform}
Das Programm sollte als eigenständige Anwendung zum Einsatz kommen. Da bisher
keine Sammlung von Kalkulationspositionen vorhanden ist, muss eine Datenbank
eingerichtet werden. Der Einsatz des My\acs{SQL}-Fork MariaDB bewährt sich
täglich im Unternehmen. \acs{PHP} dagegen ist eine Skriptsprache, die
besonders geeignet ist, um eine dynamische Webanwendung auf Basis freier
Software umzusetzen. Sie ist die am häufigsten verwendete Programmiersprache zum
Erstellen von Websites.\citet{PHP} Die Mitarbeiter sind mit den Technologien
vertraut und können gegebenfalls das System warten und weiterentwickeln.

\subsection{Architekturdesign}
\label{sec:Architekturdesign}
Der Kalkulationskatalog soll auf Basis einer \acs{MVC}-Architektur umgesetzt
werden. Diese teilt jede Softwarekomponente einen der drei
Bestandteile -- Datenmodell (engl. model), Präsentation (engl. view) oder
Programmsteuerung (engl. controller) \citet{ModelViewController} -- zu.
Jeder Teil hat einen speziellen Aufgabenbereich, der von denen der anderen
weitestgehend unabhängig ist. 

Das Model repräsentiert die Daten, der View stellt die Daten dar und der
Controller steuert den Ablauf der Anwendung -- das Bindeglied zwischen Model und
View. Durch die lose Verbindung der Komponenten, erhöht sich die
Wiederverwendbarkeit und Austauschbarkeit. So lässt sich beispielsweise die
Oberfläche komplett verändern, ohne dabei das Model anpassen zu müssen. Die
Trennung bevorteilt außerdem das Testen und Warten der Komponenten.

Es wurde sich bewusst für die Erstellung eines eigenen kleinen Frameworks
entschieden. Zum einen, weil der Lerneffekt dadurch am größten ist, zum anderen,
weil jedes vorhandene Framework eine Unmenge an nicht genutzen Code mitgebracht
hätte. Durch ein eigens entwickeltes Framework bleibt der Code schlank und
übersichtlich.

\subsection{Entwurf der Benutzeroberfläche}
\label{sec:Benutzeroberflaeche}
Eine strukturierte Oberfläche ist für eine einfache und zielgerichtete Bedienung
von großer Bedeutung. Erste Wireframes wurden mit dem Tool
Balsamiq \citet{Balsamiq} erstellt und geben die Richtung für den Aufbau der
Overfläche vor. Die Entwürfe befinden sich im \Anhang{app:Entwuerfe}.

Um schnell eine saubere,
responsive\footnote{\url{https://de.wikipedia.org/wiki/Responsive_Design}}
Weboberfläche zu generieren, wurde auf das Framework Bootstrap \citet{Bootstrap}
(ehemals \textit{Twitter Bootstrap}) gesetzt. Als \acs{CSS}-Framework bietet es
Vorlagen für alle erdenklichen Module: Textstyling, Formulare, Listen und Tabellen sind
einige wenige Vorlagen, die in diesem Projekt Verwendung finden.


\subsection{Datenmodell}
\label{sec:Datenmodell}
Das Datenmodell mit eingezeichneten Kardinalitäten ist im
\Anhang{app:Datenbankmodell} einzusehen. Im Zentrum stehen die Kalkulationen,
die jeweils einen Kunden und einen Status besitzen. Kalkulationen bestehen
außerdem aus einer Vielzahl von Positionen. Diese wiederum besitzen eine
Kategorie. Die n:m-Beziehung zwischen Kalkulationen und Positionen wird
durch eine separate Tabelle abgebildet. Globale Einstellungen, wie
Firmenname und Mehrwertsteuersatz, werden in einem gesonderten Tabelle
(\texttt{Company}) abgelegt.

\subsection{Geschäftslogik}
\label{sec:Geschaeftslogik}
Jedes Objekt innerhalb der Anwendung ist eine Instanz einer Model-Klasse. Daraus
ergeben sich folgende Klassen: \textbf{\texttt{Calculation}},
\textbf{\texttt{Category}}, \textbf{\texttt{Company}},
\textbf{\texttt{Customer}}, \textbf{\texttt{Item}} und \textbf{\texttt{Status}}.
Mit der Ausnahme von \textbf{\texttt{Company}} bekommt jede der Klassen eine
Methode \textbf{\texttt{show()}} (Anzeigen) und \textbf{\texttt{create()}}
(Anlegen), sowie \textbf{\texttt{delete()}} (Löschen), \textbf{\texttt{edit()}}
(Bearbeiten) und ggf. \textbf{\texttt{all()}} (Listenausgabe).
Weiterhin gibt es Methoden, um an Informationen anderer Tabellen zu gelangen,
\textbf{\texttt{Calculation::\-getCustomer()}} \zB liefert das Objekt des
Kunden, der der aktuellen Kalkulation zugeordnet ist. Analog dazu sind alle
anderen Methoden aufgebaut.

\subsection{Maßnahmen zur Qualitätssicherung}
\label{sec:Qualitaetssicherung}
Um die Funktionalität der Anwendung zu gewährleisten, wurden Anwendertests
geschaffen. Diese sehen die systematischen Tests aller Funktionen vor, indem
alle Aktionen der Software, mit unterschiedlichen Eingaben ausgeführt
werden.

Wird ein Test nicht bestanden, kann in git ein Fehlerbericht angelegt
werden.\footnote{\url{https://github.com/darthnorman/calc_catalog/issues}} Git
wird außerdem benutzt, um jederzeit einen stabilen Entwicklungsstrang der
Anwendung zu haben, jeder kann sich an der Weiterentwicklung beteiligen und
Vorschläge zur Verbesserung einbringen.

% !TEX root = ../Projektdokumentation.tex
\section{Implementierungsphase} 
\label{sec:Implementierungsphase}

\subsection{Implementierung der Datenstrukturen}
\label{sec:ImplementierungDatenstrukturen}
Die Datenbank wurde von Hand in Form einer \acs{SQL}-Datei angelegt. Sie wird
per phpMyAdmin in das Datenbanksystem importiert und enthält für jeden Bereich
eine kleine Anzahl an Beispieldaten. Ein Auszug der Datei liegt im
\Anhang{app:SQL}.


\subsection{Implementierung der Benutzeroberfläche}
\label{sec:ImplementierungBenutzeroberflaeche}
Bei der Implementierung der Benutzeroberfläche wurde darauf geachtet, dass eine
Übersichtliche Struktur bewart bleibt. Farben sollten sehr sparsam und nur an
Stellen eingesetzt werden, wo sie für den Nutzer eine sinnvolle Hilfestellung
darstellen. Dies macht sich vor allem bei der Darstellung von Buttons
bemerkbar. Sie stellen die primären Interaktionspunkte für Anwender dar.

Die Navigation, die es ermöglicht zwischen den Listenansichten der Komponenten
zu wechseln, ist ständig prominent am oberen Bildschirmrand platziert. Von dort
aus kann jederzeit auf eine andere Komponente der Anwendung zugegriffen werden.
Die Einstellungen (Firmenname und Mehrwertsteuersatz) werden nur extrem selten
gebraucht und wurden deshalb an den rechten Rand der Navigation gesetzt.

Ein großes Suchfeld über jeder Listenansicht bekommt bei einem Neuladen der
Seite immer den Fokus, sodass direkt mit der Eingabe und damit einer Filterung
begonnen werden kann, ohne dabei mit der Maus erst ins das Feld navigieren zu
müssen. Das Leeren des Feldes funktioniert über das Drücken der
\fbox{ESC}-Taste. Damit ist die Filterfunktion äußerst schnell rein über die
Tastatur bedienbar. Durch zusätzliche Platzhalter in den Eingabefeldern, wird
der Nutzer auf das Format seiner Eingaben hingewiesen.

Screenshots der Anwendung in der Entwicklungsphase mit Dummy-Daten befinden sich
im \Anhang{Screenshots}.


\subsection{Implementierung der Geschäftslogik}
\label{sec:ImplementierungGeschaeftslogik}
Um die Logik den Anforderungen entsprechend abzubilden, wurde zunächst eine
Konfigurationsdatei angelegt, die die Zugangsdaten für die Datenbank enthält.
Die \textit{PHP Data Objects}-Erweiterung stellt die zentrale Schnittstelle
zur Datenbank bereit und wird für alle Interaktionen mit ihr genutzt.

Für die Ausgabe wurde im nächsten Schritt die Funktion \texttt{render()}
geschrieben, die anhand der aufgerufenden \acs{URL} den Controller und den View
bestimmt. Ein Controller-Array beinhaltet alle erlaubten Kombinationen und
vergleicht diese mit der übergebenen URL, dargestellt in
\ref{fig:getCompletePriceMin}: \nameref{fig:getCompletePriceMin}.

Die einzelnen Controller bestimmen mittels des gewählten Views, welche Daten
geholt und durch welches Template diese ausgegeben werden. Das Auslesen der
Daten wird durch die Model-Klassen realisiert, indem Datenbankabfragen getätigt
werden und die Objekte an das Template übergeben werden.

Nach und nach wurden die Methoden zur Bearbeitung, Speicherung und Erstellung
neuer Datensätze geschaffen. Helferfunktionen, die kleine Aufgaben bei der
Darstellung von Werten im Template oder innerhalb von Methoden übernehmen,
wurden gesondert in der Datei \texttt{functions.php} gelistet und global
eingebunden.

Am Beispiel der Klasse \texttt{Calculation} soll die Methode
\texttt{getCompletePriceMin()} vorgestellt werden. Die Methode gibt den
minimalen Gesamtpreis (ohne Steuern) der übergebenen Kalkulation zurück.

\subsection*{\texttt{getCompletePriceMin()}-Methode}
\label{fig:getCompletePriceMin}
\begin{figure}[htb]
\begin{lstlisting}
/**
 * Berechnet den minimalen Gesamtpreis einer Kalkulation
 * @param $id ID der aktuellen Kalkulation
 * @return $pt minimaler Gesamtpreis in Euro
 */
public static function getCompletePriceMin($id) {
	// hole alle Positionen dieser Kalkulation
	$items = self::getItems($id);
	// hole aktuelle Kalkulation
	$currentCalculation = self::show($id);
	
	// Iteriere ueber alle Positionen und berechne
	// den minimalen Gesamtaufwand 
	$pt = 0;
	for ($i = 0; $i < count($items); $i++) {
		$task = $items[$i]->tmin;
		$pt = $pt + $task;
	};
	
	// Tagessatz des Teams
	$team = floatval($currentCalculation->price_team);
	
	// Multipliziere 10% das Gesamtaufwands
	// mit dem Tagessatz der PMs
	$pm = floatval($currentCalculation->price_pm);
	$pm = $pt * 0.1 * $pm;
	
	// Multipliziere die Gesamtzahl der Personentage
	// mit dem Tagessatz des Teams
	$pt = $pt * $team;
	
	// Addiere PM-Aufwand auf Gesamtaufwand
	$pt = $pt + $pm;
	
	// Gib minimalen Gesamtpreis der Kalkulation zurueck
	return $pt;
}
\end{lstlisting}
\caption{\texttt{getCompletePriceMin()}-Methode}
\end{figure}

Im Template \texttt{showCalculation.php} wird diese Methode wie folgt
aufgerufen.
\begin{lstlisting}
<?php echo getBrutto($calculation->getCompletePriceMin($calculation->id)) ?>
\end{lstlisting}
Die Funktion \texttt{getBrutto()} ist eine Hilfsfunktion und daher nicht Teil
der Klasse, sie holt sich den Mehrwertsteuersatz aus den
Einstellungen und berechnet dazu einen Wert inklusive der Steuer.
% !TEX root = ../Projektdokumentation.tex
\section{Testphase} 
\label{sec:Testphase}
Aus Zeitgründen wurden keinerlei automatische Tests implementiert. Das Testen
der Anwendung wird manuell vorgenommen, wie unter
\ref{sec:Qualitaetssicherung}: \nameref{sec:Qualitaetssicherung} beschrieben.

% !TEX root = ../Projektdokumentation.tex
\section{Einführungsphase}
\label{sec:Einfuehrungsphase}
Um die Software zu deployen, muss lediglich ein Webserver mit
einer Datenbank bereitgestellt werden. Damit eine Verbindung zur Datenbank
aufgebaut werden kann, muss die Datei \texttt{/Classes/config\_sample.php} in
\texttt{config.php} umbenannt werden und die Verbindungsdaten
(Datenbankhost, -Name, -Nutzer und Passwort) müssen eingetragen sein.
Ein Import der beiliegenden Datenbank ist nciht zwingend erforderlich, die
Datensätze können auch aus der Anwendung heraus angelegt werden.

Eine Migration alter Kalkulationen hat nicht stattgefunden, da es nicht
Bestandteil des Projektes war. Dennoch sollte der Import einiger Bestandsdaten
in naher Zukunft durchgeführt werden, es würde den Einstieg und die initiale
Arbeit mit dem Katalog erheblich erleichtern.

% !TEX root = ../Projektdokumentation.tex
\section{Dokumentation}
\label{sec:Dokumentation}

\begin{itemize}
	\item Wie wurde die Anwendung für die Benutzer/Administratoren/Entwickler dokumentiert (\zB Benutzerhandbuch, \acs{API}-Dokumentation)?
	\item Hinweis: Je nach Zielgruppe gelten bestimmte Anforderungen für die Dokumentation (\zB keine IT-Fachbegriffe in einer Anwenderdokumentation verwenden, aber auf jeden Fall in einer Dokumentation für den IT-Bereich).
\end{itemize}

\paragraph{Beispiel}
Ein Ausschnitt aus der erstellten Benutzerdokumentation befindet sich im \Anhang{app:BenutzerDoku}.
Die Entwicklerdokumentation wurde mittels PHPDoc\footnote{Vgl. \cite{phpDoc}} automatisch generiert. Ein beispielhafter Auszug aus der Dokumentation einer Klasse findet sich im \Anhang{app:Doc}. 

\Zwischenstand{Dokumentation}{Dokumentation}

% !TEX root = ../Projektdokumentation.tex
\section{Fazit} 
\label{sec:Fazit}

\subsection{Soll-/Ist-Vergleich}
\label{sec:SollIstVergleich}

\begin{itemize}
	\item Wurde das Projektziel erreicht und wenn nein, warum nicht?
	\item Ist der Auftraggeber mit dem Projektergebnis zufrieden und wenn nein, warum nicht?
	
	\item Wurde die Projektplanung (Zeit, Kosten, Personal, Sachmittel) eingehalten oder haben sich Abweichungen ergeben und wenn ja, warum?
	\item Hinweis: Die Projektplanung muss nicht strikt eingehalten werden. Vielmehr sind Abweichungen sogar als normal anzusehen. Sie müssen nur vernünftig begründet werden (\zB durch Änderungen an den Anforderungen, unter-/überschätzter Aufwand).
\end{itemize}


\subsection{Lessons Learned}
\label{sec:LessonsLearned}
Während der Bearbeitung könnte der Autor wertvolle Erfahrungen in
verschiedenen Bereichen der Softwareentwicklung machen. Zwar war mit dem
Schreiben eines eigenen \acs{MVC}-Frameworks ein nicht unerheblicher Zeitaufwand
verbunden, jedoch war der Lerneffekt dabei umso größer. Das Wissen \bzgl den
Umgang mit \LaTeX, und phpDocumentor konnte ebenfalls ausgebaut und gefestigt
werden.


\subsection{Ausblick}
\label{sec:Ausblick}
Der Katalog läuft stabil und kann im Projektalltag eingesetzt werden.
Erweiterungen sind natürlich trotzdem wünschenswert. Ein wichtiger Punkt für die
Zukunft wäre die in \namref{sec:Einfuehrungsphase} bereits erwähnte Migration
bestehnder Daten, um den Einstieg zu erleichtern.

Darüberhinaus gibt es \zB derzeit keine Möglichkeit, Positionen für eine
bestimmte Kalkulation anzupassen. Darüber hinaus könnten weitere Felder wie \bspw
Multiplikatoren für Positionen interessant sein. Auch eine Einstellmöglichkeit
für den Prozentsatz von Projektmanagerleistungen, die auf eine Kalkulation
aufgeschlagen werden, fehlt bislang. Gleiches gilt für eine Rabattfunktion.
Die modulare Struktur des Programms steht diesen Neuerungen nicht im Weg.

Wenn möglich sollten nach Abschluss eines Projektes die kalkulierten Aufgaben
mit den tatsächlich benötigten Zeiten verglichen werden. Bei Abweichungen können
die geschätzten Zeiten im Programm angepasst werden, damit zukünftige
Schätzungen noch genauer und zuverlässiger werden.

% Literatur ------------------------------------------------------------------
\clearpage
\renewcommand{\refname}{Literaturverzeichnis}
\bibliography{Bibliographie}
\bibliographystyle{Allgemein/natdin} % DIN-Stil des Literaturverzeichnisses
% !TEX root = Projektdokumentation.tex
\clearpage
\addsec{Eidesstattliche Erklärung}

% Hinweis: die eidesstattliche Erklärung ist ggfs. an die Richtlinie der IHK anzupassen

Ich, \autorName, versichere hiermit, dass ich meine
\textbf{\betreff} mit dem Thema
\begin{quote}
\textit{\kompletterTitel}
\end{quote}
selbständig verfasst und keine anderen als die angegebenen Quellen und Hilfsmittel benutzt habe,
wobei ich alle wörtlichen und sinngemäßen Zitate als solche gekennzeichnet habe. Die Arbeit
wurde bisher keiner anderen Prüfungsbehörde vorgelegt und auch nicht veröffentlicht.\\[6ex]

\abgabeOrt, den \abgabeTermin


\rule[-0.2cm]{5.5cm}{0.5pt}

\textsc{\autorName}


% Anhang ---------------------------------------------------------------------
\clearpage
\appendix
\pagenumbering{roman}
% !TEX root = Projektdokumentation.tex
\section{Anhang}
\subsection{Detaillierte Zeitplanung}
\label{app:Zeitplanung}
\tabelleAnhang{ZeitplanungKomplett}
\clearpage

\subsection{Verwendete Ressourcen}
\label{app:Verwendete Ressourcen}

\subsection*{Hardware}
\begin{itemize}
	\item Büroarbeitsplatz mit Desktop-PC
\end{itemize}

\subsection*{Software}
\begin{itemize}
	\item Windows 10 Pro 64\,bit mit Service Pack 1 -- Betriebssystem
	\item Eclipse Mars mit TeXlipse -- Entwicklungsumgebung mit Erweiterung für
	\LaTeX
	\item \acs{XAMPP} -- Apache-Distribution, die MariaDB, PHP und Perl enthält
	\item MariaDB -- Datenbanksystem
	\item phpMyAdmin -- Datenbankmanagementsystem
	\item phpDocumentor -- Software zur Erstellung von Entwicklerdokumentationen
	\item git -- dezentrale Versionsverwaltung
	\item MiKTeX -- Distribution des Textsatzsystems \TeX
	\item Gimp -- Bildbearbeitungsprogramm
	\item Bootstrap -- Framework für Oberflächen 
	\item jQuery -- JavaScript-Bibliothek
	\item List.js -- Bibliothek zum Sortieren und Filtern von Listen im \acs{FE}
	\item Balsamiq -- Programm zur Erstellung von Wireframes
	\item Mozilla Firefox -- Internetbrowser zum Betrachten von Webanwendungen
\end{itemize}

\subsection*{Personal}
\begin{itemize}
	\item Projektmanagerin -- Festlegung der Anforderungen und Abnahme  
	\item Entwickler -- Umsetzung des Projektes
	\item Anwendungsentwickler -- Beratung
\end{itemize}
\clearpage

\subsection{Ausschnitt einer Excel-Tabelle}
\label{app:Exceltabelle}
\begin{figure}[htb]
\centering
\includegraphicsKeepAspectRatio{istanalyse.png}{1}
\caption{Ausschnitt einer Excel-Tabelle einer Kalkulation}
\end{figure}
\clearpage

\subsection{Datenbankmodell}
\label{app:Datenbankmodell}
\begin{figure}[htb]
\centering
\includegraphicsKeepAspectRatio{datenbankschema.png}{1}
\caption{Datenbankmodell}
\end{figure}
\clearpage

\subsection{SQL - Auszug}
\label{app:SQL}
\begin{figure}[htb]
\begin{lstlisting}
CREATE TABLE `item` (
  `id` int(6) unsigned NOT NULL auto_increment,
  `name` varchar(128) collate utf8_unicode_ci NOT NULL,
  `description` text(1000),
  `tmin` float(6,3) unsigned NOT NULL DEFAULT '0',
  `tmax` float(6,3) unsigned NOT NULL DEFAULT '0',
  `category` int(6) unsigned NOT NULL,
  PRIMARY KEY (`id`),
  CONSTRAINT `fk_category` FOREIGN KEY (`category`) REFERENCES category(`id`)
) ENGINE=InnoDB  DEFAULT CHARSET=utf8 COLLATE=utf8_unicode_ci AUTO_INCREMENT=8;

INSERT INTO `item` VALUES(1, 'Berechtigungen vergeben', 'Den Nutzern werden
Rechte vergeben, damit sie ihre Aufgaben erledigen koennen.', 1.25, 2, 1);
\end{lstlisting}
\caption{Auszug aus der Datenbank in SQL-Format}
\end{figure}
\clearpage

\subsection{Oberflächenentwürfe}
\label{app:Entwuerfe}
\begin{figure}[htb]
\centering
\includegraphicsKeepAspectRatio{MockupModules.pdf}{0.7}
\caption{Liste der Module mit Filtermöglichkeiten}
\end{figure}

\begin{figure}[htb]
\centering
\includegraphicsKeepAspectRatio{MockupModul.pdf}{0.7}
\caption{Anzeige der Übersichtsseite einzelner Module}
\end{figure}

\begin{figure}[htb]
\centering
\includegraphicsKeepAspectRatio{MockupTag.pdf}{0.7}
\caption{Anzeige und Filterung der Module nach Tags}
\end{figure}

\clearpage
\subsection{Screenshots der Anwendung}
\label{Screenshots}
\begin{figure}[htb]
\centering
\includegraphicsKeepAspectRatio{tagliste.pdf}{1}
\caption{Anzeige und Filterung der Module nach Tags}
\end{figure}
\clearpage
\begin{figure}[htb]
\centering
\includegraphicsKeepAspectRatio{modulliste.pdf}{1}
\caption{Liste der Module mit Filtermöglichkeiten}
\end{figure}
\clearpage

\clearpage


\end{document}
